\section{Deferred proofs}
\label{sec:app}

% \begin{proof}[Proof of Lemma~\ref{lemma:master}]
%     Dividing numerator and denominator of the last fraction
%     in~\eqref{eq:Lik} by $g^{n/2}$, and writing $C_n(Y)$ for the
%     $\gamma$-free prefactor,
%     \begin{equation}
%         \label{eq:logL}
%         \log\mathcal{L}_n(Y\mid\gamma)
%         =\log C_n(Y)-\frac n2\log g
%         -\frac{\abs \gamma}{2}\log(1+g)
%         -\frac n2\log\bigl\{g^{-1}+1-R^{2}_\gamma\bigr\}.
%     \end{equation}
%     The first two terms of the right-hand side of~\eqref{eq:logL} do not
%     depend on $\gamma$, hence cancel in a difference:
%     \begin{equation}
%         \label{eq:logLratio}
%         \log\frac{\mathcal{L}_n(Y\mid\gamma)}
%         {\mathcal{L}_n(Y\mid\gamma')}
%         =\frac{\abs{\gamma'}-\abs \gamma}{2}\log(1+g)
%         +\frac n2\log\frac{g^{-1}+1-R^{2}_{\gamma'}}
%         {g^{-1}+1-R^{2}_{\gamma}} .
%     \end{equation}
%     By~\eqref{eq:prior} and~\eqref{eq:pilambda}, the normalising constants
%     cancelling,
%     \begin{equation*}
%         \log\frac{\pi_n^{(\lambda)}(\gamma\mid Y)}
%         {\pi_n^{(\lambda)}(\gamma'\mid Y)}
%         =\log\frac{\pi(\gamma)}{\pi(\gamma')}
%         +\lambda\log\frac{\mathcal{L}_n(Y\mid\gamma)}
%         {\mathcal{L}_n(Y\mid\gamma')}
%         =\kappa\bigl(\abs{\gamma'}-\abs \gamma\bigr)\log p
%         +\lambda\log\frac{\mathcal{L}_n(Y\mid\gamma)}
%         {\mathcal{L}_n(Y\mid\gamma')} .
%     \end{equation*}
%     Substituting~\eqref{eq:logLratio} and collecting the two terms
%     carrying the factor $\abs{\gamma'}-\abs \gamma$
%     gives~\eqref{eq:master}.
% \end{proof}

Every proof of the paper is collected here, together with the notation and
the deterministic bounds they rest on. Only the proof of \cref{th:drop}
involves the temperature, and it involves it only through the factor
$\lambda$ multiplying the fit term of~\eqref{eq:master}.

\subsection{Notation}
\label{sec:prelim}

We use throughout the effective noise of YWJ,
\begin{equation}
    \label{eq:wtilde}
    \widetilde w\ \coloneq\ X_{S^{c}}\beta^\ast_{S^{c}}+w,
    \qquad\text{so that}\qquad
    Y=\signal+\widetilde w ,
\end{equation}
the elementary inequalities
\begin{gather}
    \label{eq:elementary}
    \log(1+x)\leq x\ \ (x>-1),
    \qquad
    \norm{a+b}_2^{2}\leq2\norm a_2^{2}+2\norm b_2^{2},\\
    \nonumber
    \frac{a}{a+b}\geq\min\Bigl\{\frac12,\frac{a}{2b}\Bigr\}
    \ \ (a,b>0),
\end{gather}
and, for an underfitted $\gamma\in\mathscr{M}(s_0)$, the quantity
\begin{equation}
    \label{eq:Adef}
    A^{2}\ \coloneq\ \nu\,
    \frac{\norm{(I-\Phi_\gamma)\signal}_2^{2}}
    {\abs{\gamma^{\ast}\setminus\gamma}}\ >\ 0 ,
\end{equation}
the dependence on $\gamma$ being suppressed, as YWJ suppress it, since
$\gamma$ is fixed wherever $A$ occurs.


\subsection{Deterministic bounds}
\label{sub:appdet}

\begin{lemma}[Uniform residual lower bound]
    \label{lemma:residual}
    On the good event, and under~\eqref{eq:Dassumption}, every
    $\bar\gamma\in\mathscr{M}(s_0)$ satisfies
    \begin{equation}
        \label{eq:residual}
        \norm Y_2^{2}\bigl(g^{-1}+1-R^{2}_{\bar\gamma}\bigr)
        \ \geq\ \norm{(I-\Phi_{\bar\gamma})Y}_2^{2}
        \ \geq\ \frac{n\sigma_0^{2}}{8} .
    \end{equation}
\end{lemma}

\begin{proof}
    Since $\norm Y_2^{2}(g^{-1}+1-R^{2}_{\bar\gamma})
    =g^{-1}\norm Y_2^{2}+\norm{(I-\Phi_{\bar\gamma})Y}_2^{2}$, the first
    inequality amounts to discarding $g^{-1}\norm Y_2^{2}\geq0$. For the
    second, fix $\bar\gamma\in\mathscr{M}(s_0)$ and set
    $\bar\gamma^{+}\coloneq\bar\gamma\cup\gamma^{\ast}$, so that
    $\abs{\bar\gamma^{+}}\leq\abs{\bar\gamma}+s^{\ast}\leq s$ and
    $\bar\gamma^{+}\in\mathscr{M}(s)$. Since
    $\bar\gamma\subseteq\bar\gamma^{+}$, the range of $\Phi_{\bar\gamma}$
    is contained in that of $\Phi_{\bar\gamma^{+}}$, so
    $I-\Phi_{\bar\gamma}\succcurlyeq I-\Phi_{\bar\gamma^{+}}$ and
    \begin{equation}
        \label{eq:res1}
        \norm{(I-\Phi_{\bar\gamma})Y}_2^{2}
        \ \geq\ \norm{(I-\Phi_{\bar\gamma^{+}})Y}_2^{2}.
    \end{equation}
    Because $\gamma^{\ast}\subseteq\bar\gamma^{+}$, the vector $\signal$
    lies in the range of $\Phi_{\bar\gamma^{+}}$, so
    $(I-\Phi_{\bar\gamma^{+}})\signal=0$ and
    \begin{equation*}
    (I-\Phi_{\bar\gamma^{+}})
        Y
        =(I-\Phi_{\bar\gamma^{+}})w
        +(I-\Phi_{\bar\gamma^{+}})X_{S^{c}}\beta^\ast_{S^{c}} .
    \end{equation*}
    Apply $\norm{a+b}_2^{2}\geq\tfrac12\norm a_2^{2}-\norm b_2^{2}$, which
    is the rearrangement of $\tfrac12(\norm a_2-2\norm b_2)^{2}\geq0$,
    with $a=(I-\Phi_{\bar\gamma^{+}})w$ and
    $b=(I-\Phi_{\bar\gamma^{+}})X_{S^{c}}\beta^\ast_{S^{c}}$. Since
    projections are non-expansive, $\norm b_2^{2}\leq
    \norm{X_{S^{c}}\beta^\ast_{S^{c}}}_2^{2}\leq\Ltil\sigma_0^{2}\log p$
    by \cref{ass:A}, and $\norm a_2^{2}=\norm w_2^{2}
    -w^{\top}\Phi_{\bar\gamma^{+}}w$ because $I-\Phi_{\bar\gamma^+}$ is a
    projection. Hence
    \begin{equation*}
        \norm{(I-\Phi_{\bar\gamma^{+}})Y}_2^{2}
        \ \geq\ \tfrac12\bigl(\norm w_2^{2}
        -w^{\top}\Phi_{\bar\gamma^{+}}w\bigr)
        -\Ltil\sigma_0^{2}\log p .
    \end{equation*}
    On $\mathcal{C}_n$ we have $\norm w_2^{2}\geq n\sigma_0^{2}/2$, and on
    $\mathcal{B}_n$, applied to $\bar\gamma^{+}\in\mathscr{M}(s)$, we have
    $w^{\top}\Phi_{\bar\gamma^{+}}w\leq 8s\,\sigma_0^{2}\log p$.
    Therefore
    \begin{equation*}
        \norm{(I-\Phi_{\bar\gamma^{+}})Y}_2^{2}
        \ \geq\ \frac{n\sigma_0^{2}}{4}
        -\bigl\{4s+\Ltil\bigr\}\sigma_0^{2}\log p
        \ \geq\ \frac{n\sigma_0^{2}}{4}
        -\bigl\{8s+\Ltil\bigr\}\sigma_0^{2}\log p ,
    \end{equation*}
    and by~\eqref{eq:Dassumption} the subtracted term is at most
    $n\sigma_0^{2}/32$, so the right-hand side is at least
    $7n\sigma_0^{2}/32\geq n\sigma_0^{2}/8$. Combining
    with~\eqref{eq:res1} finishes the proof.
\end{proof}

\begin{lemma}[Effective noise]
    \label{lemma:noise}
    On the good event, and under~\eqref{eq:Dassumption}:
    \begin{equivcond}
        \item\label{it:noise_i}
        $\norm{\widetilde w}_2^{2}\leq2\Ltil\sigma_0^{2}\log p
        +3n\sigma_0^{2}$;
        \item\label{it:noise_ii} for every pair
        $\gamma_2\subset\gamma_1$ of $\mathscr{M}(s)$ with
        $\abs{\gamma_1}-\abs{\gamma_2}=1$,
        \begin{equation*}
            \norm{(\Phi_{\gamma_1}-\Phi_{\gamma_2})\widetilde w}_2^{2}
            \ \leq\ 2(L+\Ltil)\,\sigma_0^{2}\log p
            \ \leq\ \frac{n\nu^{2}C_\beta^{2}}{64} .
        \end{equation*}
    \end{equivcond}
\end{lemma}

\begin{proof}
    \ref{it:noise_i} By~\eqref{eq:elementary} and \cref{ass:A},
    $\norm{\widetilde w}_2^{2}\leq2\norm{X_{S^{c}}
        \beta^\ast_{S^{c}}}_2^{2}+2\norm w_2^{2}
    \leq2\Ltil\sigma_0^{2}\log p+2\norm w_2^{2}$, and on $\mathcal{C}_n$
    we have $\norm w_2^{2}\leq\tfrac32n\sigma_0^{2}$.

    \ref{it:noise_ii} Write $\Pi=\Phi_{\gamma_1}-\Phi_{\gamma_2}$, an
    orthogonal projection, so that $\norm{\Pi v}_2\leq\norm v_2$ for every
    $v$ and $\norm{\Pi w}_2^{2}=w^{\top}\Pi w$. Hence,
    by~\eqref{eq:elementary},
    \begin{align*}
        \norm{\Pi\widetilde w}_2^{2}
        &\leq2\norm{\Pi X_{S^{c}}\beta^\ast_{S^{c}}}_2^{2}
        +2\norm{\Pi w}_2^{2}
        \leq2\norm{X_{S^{c}}\beta^\ast_{S^{c}}}_2^{2}
        +2\,w^{\top}\bigl(\Phi_{\gamma_1}-\Phi_{\gamma_2}\bigr)w\\
        &\leq2\Ltil\sigma_0^{2}\log p+2L\sigma_0^{2}\log p,
    \end{align*}
    using \cref{ass:A} for the first term and the event $\mathcal{A}_n$,
    applied to the nested pair $\gamma_2\subset\gamma_1$ of
    $\mathscr{M}(s)$ with $\abs{\gamma_1}-\abs{\gamma_2}=1$, for the
    second. The final inequality is \cref{ass:betamin} rearranged.
\end{proof}

\begin{lemma}[Upper bound on the regularised residual]
    \label{lemma:Dupper}
    On the good event, and under~\eqref{eq:Dassumption} and $n\geq3$, every
    $\gamma\in\mathscr{M}(s_0)$ satisfies
    \begin{equation}
        \label{eq:Dupper}
        \norm Y_2^{2}\bigl(g^{-1}+1-R^{2}_{\gamma}\bigr)
        \ \leq\ 2\norm{(I-\Phi_\gamma)\signal}_2^{2}
        +8\,n\sigma_0^{2} .
    \end{equation}
\end{lemma}

\begin{proof}
    By~\eqref{eq:wtilde} and~\eqref{eq:elementary},
    \begin{equation*}
        \norm{(I-\Phi_\gamma)Y}_2^{2}
        \leq2\norm{(I-\Phi_\gamma)\signal}_2^{2}
        +2\norm{(I-\Phi_\gamma)\widetilde w}_2^{2}
        \leq2\norm{(I-\Phi_\gamma)\signal}_2^{2}
        +2\norm{\widetilde w}_2^{2},
    \end{equation*}
    the last step because $I-\Phi_\gamma$ is a projection. By
    \cref{lemma:noise}\ref{it:noise_i}, $2\norm{\widetilde w}_2^{2}\leq
    4\Ltil\sigma_0^{2}\log p+6n\sigma_0^{2}$, and on $\mathcal{D}_n$ we
    have $g^{-1}\norm Y_2^{2}\leq(2\log p+3)\sigma_0^{2}$. Adding the two,
    the additive term is at most
    \begin{equation*}
        \bigl\{4\Ltil\log p+6n+2\log p+3\bigr\}\sigma_0^{2}
        \ \leq\ \Bigl\{\frac n8+6n+\frac{n}{128}+n\Bigr\}\sigma_0^{2}
        \ \leq\ 8n\sigma_0^{2},
    \end{equation*}
    where we used, in order: $\Ltil\log p\leq n/32$
    from~\eqref{eq:Dassumption}, whence $4\Ltil\log p\leq n/8$; the same display
    with $\Ltil$ replaced by $8s\geq8$, whence $\log p\leq n/256$ and
    $2\log p\leq n/128$; and $3\leq n$.
\end{proof}

\begin{lemma}[Greedy selection; YWJ, Lemma~8]
    \label{lemma:forward}
    Let $\gamma\in\mathscr{M}(s_0)$ be underfitted and let $j_\gamma$ be
    as in~\eqref{eq:jdef}. Under \cref{ass:B,ass:D,ass:betamin}:
    \begin{equivcond}
        \item\label{it:fs_a}
        $\norm{(\Phi_{\gamma\cup\{j_\gamma\}}-\Phi_\gamma)
            \signal}_2^{2}\geq A^{2}$, and moreover
        \begin{equation}
            \label{eq:Alower}
            A^{2}\ \geq\ n\,\nu^{2}\min_{j\in\gamma^{\ast}}
            \abs{\beta^\ast_j}^{2}
            \ \geq\ n\,\nu^{2}C_\beta^{2} ;
        \end{equation}
        \item\label{it:fs_b} if in addition $\gamma$ is saturated and
        $k_\gamma$ is as in~\eqref{eq:kdef}, then, with
        $\widetilde\gamma=\gamma\cup\{j_\gamma\}$ and
        $\gamma'=\widetilde\gamma\setminus\{k_\gamma\}$,
        \begin{equation*}
            \norm{(\Phi_{\widetilde\gamma}-\Phi_{\gamma'})
                \signal}_2^{2}
            \ \leq\ n\,\omega(X)\,
            \frac{\norm{\beta^\ast_{\gamma^{\ast}\setminus\gamma}}_2^{2}}
            {s_0-s^{\ast}}
            \ \leq\ \frac{A^{2}}{4} .
        \end{equation*}
    \end{equivcond}
\end{lemma}

\begin{proof}
    We use two facts from YWJ. The first is their Lemma~7, the rank-one
    update formula: for $\bar\gamma\in\mathscr{M}(s)$ and
    $k\notin\bar\gamma$ with $\abs{\bar\gamma}<s$,
    \begin{equation}
        \label{eq:rankone}
        \Phi_{\bar\gamma\cup\{k\}}-\Phi_{\bar\gamma}
        =\frac{(I-\Phi_{\bar\gamma})X_kX_k^{\top}
            (I-\Phi_{\bar\gamma})}
        {X_k^{\top}(I-\Phi_{\bar\gamma})X_k},
    \end{equation}
    which follows from the block inversion formula applied to
    $X_{\bar\gamma\cup\{k\}}^{\top}X_{\bar\gamma\cup\{k\}}$. The second
    is their Lemma~6: under \cref{ass:B}, for every nested pair
    $\gamma_2\subset\gamma_1$ of $\mathscr{M}(s)$,
    \begin{equation}
        \label{eq:underfit}
        \lambda_{\min}\Bigl(\frac1n
        X_{\gamma_1\setminus\gamma_2}^{\top}
        (I-\Phi_{\gamma_2})X_{\gamma_1\setminus\gamma_2}\Bigr)
        \ \geq\ \nu ,
    \end{equation}
    because the matrix inside is $n$ times the Schur complement of the
    block $X_{\gamma_2}^{\top}X_{\gamma_2}$ in
    $X_{\gamma_1}^{\top}X_{\gamma_1}$, hence its inverse is the
    corresponding diagonal block of
    $(n^{-1}X_{\gamma_1}^{\top}X_{\gamma_1})^{-1}$, whose largest
    eigenvalue is at most
    $\lambda_{\min}(n^{-1}X_{\gamma_1}^{\top}X_{\gamma_1})^{-1}
    \leq\nu^{-1}$.

    \emph{Proof of~\ref{it:fs_a}.} Put $u\coloneq(I-\Phi_\gamma)\signal$.
    For $\ell\in\gamma^{\ast}\setminus\gamma$, idempotence of projections
    and~\eqref{eq:rankone} give
    \begin{align}
        \notag
        \norm{\Phi_{\gamma\cup\{\ell\}}\signal}_2^{2}
        -\norm{\Phi_{\gamma}\signal}_2^{2}
        &=(\beta^\ast_{\gamma^{\ast}})^{\top}
        X_{\gamma^{\ast}}^{\top}
        \bigl(\Phi_{\gamma\cup\{\ell\}}-\Phi_\gamma\bigr)\signal\\
        \label{eq:onegain}
        &=\frac{\langle X_\ell,u\rangle^{2}}
        {X_\ell^{\top}(I-\Phi_\gamma)X_\ell}
        \ \geq\ \frac{\langle X_\ell,u\rangle^{2}}{n},
    \end{align}
    where the last step uses
    $X_\ell^{\top}(I-\Phi_\gamma)X_\ell\leq\norm{X_\ell}_2^{2}=n$ and,
    for the middle equality, that
    $X_\ell^{\top}(I-\Phi_\gamma)\signal=\langle X_\ell,u\rangle$.
    Summing~\eqref{eq:onegain} over $\ell\in\gamma^{\ast}\setminus\gamma$,
    \begin{equation}
        \label{eq:sumgain}
        \sum_{\ell\in\gamma^{\ast}\setminus\gamma}
        \Bigl\{\norm{\Phi_{\gamma\cup\{\ell\}}\signal}_2^{2}
        -\norm{\Phi_{\gamma}\signal}_2^{2}\Bigr\}
        \ \geq\ \frac1n
        \norm{X_{\gamma^{\ast}\setminus\gamma}^{\top}u}_2^{2}
        \ =\ \frac1n\norm{X_{\gamma^{\ast}\cup\gamma}^{\top}u}_2^{2},
    \end{equation}
    the last equality because $X_j^{\top}u=0$ for $j\in\gamma$: indeed
    $X_j$ lies in the range of $\Phi_\gamma$, to which $u$ is orthogonal
    by construction. Now $u$ lies in the range of
    $X_{\gamma^{\ast}\cup\gamma}$, so $u=X_{\gamma^{\ast}\cup\gamma}v$ for
    some $v$; writing
    $M\coloneq X_{\gamma^{\ast}\cup\gamma}^{\top}
    X_{\gamma^{\ast}\cup\gamma}$, we have
    $\abs{\gamma^{\ast}\cup\gamma}\leq s$, hence $M\succcurlyeq n\nu I$ by
    \cref{ass:B}, and therefore
    \begin{equation*}
        \frac1n\norm{X_{\gamma^{\ast}\cup\gamma}^{\top}u}_2^{2}
        =\frac1nv^{\top}M^{2}v
        \ \geq\ \nu\,v^{\top}Mv
        =\nu\norm u_2^{2},
    \end{equation*}
    the inequality because $M^{2}-n\nu M=M^{1/2}(M-n\nu I)M^{1/2}
    \succcurlyeq0$. Combining with~\eqref{eq:sumgain} and dividing by
    $\abs{\gamma^{\ast}\setminus\gamma}$, the maximum over $\ell$ of the
    left-hand side of~\eqref{eq:onegain} is at least
    $\nu\norm u_2^{2}/\abs{\gamma^{\ast}\setminus\gamma}=A^{2}$. Since
    $j_\gamma$ attains that maximum by~\eqref{eq:jdef}, and since the
    left-hand side of~\eqref{eq:onegain} equals
    $\norm{(\Phi_{\gamma\cup\{\ell\}}-\Phi_\gamma)\signal}_2^{2}$ by
    idempotence, the first claim follows.

    For~\eqref{eq:Alower}, note that
    $(I-\Phi_\gamma)\signal=(I-\Phi_\gamma)
    X_{\gamma^{\ast}\setminus\gamma}
    \beta^\ast_{\gamma^{\ast}\setminus\gamma}$, because the columns $X_j$
    with $j\in\gamma^{\ast}\cap\gamma$ lie in the range of $\Phi_\gamma$;
    hence by~\eqref{eq:underfit} applied to
    $\gamma\subset\gamma^{\ast}\cup\gamma$,
    \begin{equation}
        \label{eq:usquare}
        \norm u_2^{2}
        =(\beta^\ast_{\gamma^{\ast}\setminus\gamma})^{\top}
        X_{\gamma^{\ast}\setminus\gamma}^{\top}(I-\Phi_\gamma)
        X_{\gamma^{\ast}\setminus\gamma}
        \beta^\ast_{\gamma^{\ast}\setminus\gamma}
        \ \geq\ n\,\nu
        \norm{\beta^\ast_{\gamma^{\ast}\setminus\gamma}}_2^{2},
    \end{equation}
    so that, using
    $\norm{\beta^\ast_{\gamma^{\ast}\setminus\gamma}}_2^{2}
    \geq\abs{\gamma^{\ast}\setminus\gamma}
    \min_{j\in\gamma^{\ast}}\abs{\beta^\ast_j}^{2}$ and~\eqref{eq:S},
    \begin{equation}
        \label{eq:Achain}
        A^{2}=\frac{\nu\norm u_2^{2}}
        {\abs{\gamma^{\ast}\setminus\gamma}}
        \ \geq\ \frac{n\nu^{2}
            \norm{\beta^\ast_{\gamma^{\ast}\setminus\gamma}}_2^{2}}
        {\abs{\gamma^{\ast}\setminus\gamma}}
        \ \geq\ n\nu^{2}\min_{j\in\gamma^{\ast}}
        \abs{\beta^\ast_j}^{2}
        \ \geq\ n\nu^{2}C_\beta^{2} .
    \end{equation}

    \emph{Proof of~\ref{it:fs_b}.} Write
    $\widetilde\gamma=\gamma\cup\{j_\gamma\}$. For
    $\bar\gamma\in\mathscr{M}(s)$ let $\widehat\beta(\bar\gamma)$ solve
    \begin{equation}
        \label{eq:modLS}
        \min\bigl\{\norm{X\beta
        -X_{\gamma^{\ast}\setminus\widetilde\gamma}
            \beta^\ast_{\gamma^{\ast}\setminus\widetilde\gamma}}_2^{2}:\
        \beta\in\bbR^{p},\ \beta_j=0\ \text{for }j\notin\bar\gamma\bigr\},
    \end{equation}
    so that
    \begin{equation}
        \label{eq:LSid}
        \norm{X\widehat\beta(\bar\gamma)
        -X_{\gamma^{\ast}\setminus\widetilde\gamma}
            \beta^\ast_{\gamma^{\ast}\setminus\widetilde\gamma}}_2^{2}
        =\norm{(I-\Phi_{\bar\gamma})
            X_{\gamma^{\ast}\setminus\widetilde\gamma}
            \beta^\ast_{\gamma^{\ast}\setminus\widetilde\gamma}}_2^{2},
        \qquad
        \widehat\beta(\bar\gamma)_{\bar\gamma}
        =(X_{\bar\gamma}^{\top}X_{\bar\gamma})^{-1}X_{\bar\gamma}^{\top}
        X_{\gamma^{\ast}\setminus\widetilde\gamma}
        \beta^\ast_{\gamma^{\ast}\setminus\widetilde\gamma} .
    \end{equation}
    Since
    $k_\gamma\notin\gamma^\ast$ we have
    $\Phi_{\widetilde\gamma}X_{\gamma^{\ast}\cap\widetilde\gamma}
    =\Phi_{\widetilde\gamma\setminus\{k\}}
    X_{\gamma^{\ast}\cap\widetilde\gamma}$ for every
    $k\in\widetilde\gamma\setminus\gamma^{\ast}$, so with
    $\gamma'=\widetilde\gamma\setminus\{k_\gamma\}$
    \begin{align}
        \notag
        \norm{(\Phi_{\widetilde\gamma}-\Phi_{\gamma'})
            \signal}_2^{2}
        &=\norm{(\Phi_{\widetilde\gamma}-\Phi_{\gamma'})
            X_{\gamma^{\ast}\setminus\widetilde\gamma}
            \beta^\ast_{\gamma^{\ast}\setminus\widetilde\gamma}}_2^{2}\\
        \notag
        &=\norm{(I-\Phi_{\gamma'})
            X_{\gamma^{\ast}\setminus\widetilde\gamma}
            \beta^\ast_{\gamma^{\ast}\setminus\widetilde\gamma}}_2^{2}
        -\norm{(I-\Phi_{\widetilde\gamma})
            X_{\gamma^{\ast}\setminus\widetilde\gamma}
            \beta^\ast_{\gamma^{\ast}\setminus\widetilde\gamma}}_2^{2}\\
        \label{eq:LSdiff}
        &=\norm{X\widehat\beta(\gamma')
        -X_{\gamma^{\ast}\setminus\widetilde\gamma}
            \beta^\ast_{\gamma^{\ast}\setminus\widetilde\gamma}}_2^{2}
        -\norm{X\widehat\beta(\widetilde\gamma)
        -X_{\gamma^{\ast}\setminus\widetilde\gamma}
            \beta^\ast_{\gamma^{\ast}\setminus\widetilde\gamma}}_2^{2},
    \end{align}
    using in the second line that
    $\Phi_{\widetilde\gamma}-\Phi_{\gamma'}$ is a projection and
    $\gamma'\subset\widetilde\gamma$. By optimality of
    $\widehat\beta(\widetilde\gamma)$ in~\eqref{eq:modLS}, the residual
    $X\widehat\beta(\widetilde\gamma)
    -X_{\gamma^{\ast}\setminus\widetilde\gamma}
    \beta^\ast_{\gamma^{\ast}\setminus\widetilde\gamma}$ is orthogonal to
    $X_k$ for every $k\in\widetilde\gamma$, so for such $k$
    \begin{align*}
        \norm{X\widehat\beta(\widetilde\gamma\setminus\{k\})
        -X_{\gamma^{\ast}\setminus\widetilde\gamma}
            \beta^\ast_{\gamma^{\ast}\setminus\widetilde\gamma}}_2^{2}
        &\ \leq\ \norm{X\widehat\beta(\widetilde\gamma)
        -X_{\gamma^{\ast}\setminus\widetilde\gamma}
            \beta^\ast_{\gamma^{\ast}\setminus\widetilde\gamma}
            -X_k\widehat\beta_k(\widetilde\gamma)}_2^{2}\\
        &\ =\ \norm{X\widehat\beta(\widetilde\gamma)
        -X_{\gamma^{\ast}\setminus\widetilde\gamma}
            \beta^\ast_{\gamma^{\ast}\setminus\widetilde\gamma}}_2^{2}
        +\norm{X_k}_2^{2}\,
        \abs{\widehat\beta_k(\widetilde\gamma)}^{2},
    \end{align*}
    the first line because the vector $\widehat\beta(\widetilde\gamma)$
    with its $k$th coordinate zeroed is feasible for the problem defining
    $\widehat\beta(\widetilde\gamma\setminus\{k\})$, and the second by
    orthogonality. Since $\norm{X_k}_2^{2}=n$ and $k_\gamma$ minimises the
    left-hand side over $k\in\gamma\setminus\gamma^\ast
    =\widetilde\gamma\setminus(\gamma^{\ast}\cup\{j_\gamma\})$, whose
    cardinality is $s_0-\abs{\gamma\cap\gamma^{\ast}}\geq s_0-s^{\ast}$,
    \begin{equation*}
        \norm{X\widehat\beta(\gamma')
        -X_{\gamma^{\ast}\setminus\widetilde\gamma}
            \beta^\ast_{\gamma^{\ast}\setminus\widetilde\gamma}}_2^{2}
        -\norm{X\widehat\beta(\widetilde\gamma)
        -X_{\gamma^{\ast}\setminus\widetilde\gamma}
            \beta^\ast_{\gamma^{\ast}\setminus\widetilde\gamma}}_2^{2}
        \ \leq\ n\min_{k\in\gamma\setminus\gamma^\ast}
        \abs{\widehat\beta_k(\widetilde\gamma)}^{2}
        \ \leq\ \frac{n\norm{\widehat\beta(\widetilde\gamma)}_2^{2}}
        {s_0-s^{\ast}},
    \end{equation*}
    a minimum being at most the average. Finally
    \begin{equation*}
        \norm{\widehat\beta(\widetilde\gamma)}_2^{2}
        =\norm{(X_{\widetilde\gamma}^{\top}X_{\widetilde\gamma})^{-1}
            X_{\widetilde\gamma}^{\top}
            X_{\gamma^{\ast}\setminus\widetilde\gamma}
            \beta^\ast_{\gamma^{\ast}\setminus\widetilde\gamma}}_2^{2}
        \ \leq\ \omega(X)\,
        \norm{\beta^\ast_{\gamma^{\ast}
        \setminus\widetilde\gamma}}_2^{2}
        \ \leq\ \omega(X)\,
        \norm{\beta^\ast_{\gamma^{\ast}\setminus\gamma}}_2^{2},
    \end{equation*}
    by the definition of $\omega(X)$ in \cref{ass:D} and
    $\gamma^{\ast}\setminus\widetilde\gamma\subseteq
    \gamma^{\ast}\setminus\gamma$. Together with~\eqref{eq:LSdiff} this
    gives the first inequality of~\ref{it:fs_b}.

    For the second, combine
    $A^{2}\geq n\nu^{2}
    \norm{\beta^\ast_{\gamma^{\ast}\setminus\gamma}}_2^{2}/
    \abs{\gamma^{\ast}\setminus\gamma}\geq
    n\nu^{2}\norm{\beta^\ast_{\gamma^{\ast}\setminus\gamma}}_2^{2}/
    s^{\ast}$, which is~\eqref{eq:Achain}, with the first inequality of
    \cref{ass:D}: $s_0\geq(4\nu^{-2}\omega(X)+1)s^{\ast}$ is equivalent to
    $4\omega(X)s^{\ast}\leq\nu^{2}(s_0-s^{\ast})$, that is to
    \begin{equation*}
        \frac{n\,\omega(X)
        \norm{\beta^\ast_{\gamma^{\ast}\setminus\gamma}}_2^{2}}
        {s_0-s^{\ast}}
        \ \leq\ \frac14\cdot
        \frac{n\,\nu^{2}
            \norm{\beta^\ast_{\gamma^{\ast}\setminus\gamma}}_2^{2}}
        {s^{\ast}}
        \ \leq\ \frac{A^{2}}{4} . \qedhere
    \end{equation*}
\end{proof}


\begin{remark}
    \label{rem:ywj_constant}
    With their constant $2$ in place of our $4$ in~\eqref{eq:Dassumption},
    YWJ obtain $\norm{v_2}_2\leq A/\sqrt2$ rather than
    $\norm{v_2}_2\leq A/2$, and combine it with a noise bound $A/4$. The
    resulting numbers do not close: on their page~2527 the displayed lower
    bound $A(A-A/4)-(A/\sqrt2)(A/\sqrt2+A/4)-A^{2}/16$ equals
    $0.0107\,A^{2}$, not the $A^{2}/8$ claimed, and it is negative if one
    keeps the factor $2$ that their own preceding display carries on the
    cross terms. Strengthening \cref{ass:D} to
    $s_0\geq(4\nu^{-2}\omega(X)+1)s^{\ast}$, and \cref{ass:betamin} so
    that the noise budget is $A/8$ rather than $A/4$, repairs this with
    room to spare; see~\eqref{eq:satdrop} below. Nothing else in their
    argument, or in ours, is affected.
\end{remark}

\begin{lemma}[Proxy bound]
    \label{lemma:proxy}
    For every $\gamma\in\mathscr{M}(s)$ and every
    $j\in\gamma\setminus\gamma^{\ast}$,
    \begin{equation}
        \label{eq:proxy}
        \bigl\|(\Phi_{\gamma}-\Phi_{\gamma\setminus\{j\}})
        \signal\bigr\|_2^{2}
        \ \leq\ n\,\omega(X)\,
        \bigl\|\beta^{\ast}_{\gamma^{\ast}\setminus\gamma}\bigr\|_2^{2} .
    \end{equation}
\end{lemma}

\begin{proof}
    This is the chain of displays in the proof of
    \cref{lemma:forward}\ref{it:fs_b}, run with $\widetilde\gamma$
    replaced by $\gamma$ and $k_\gamma$ by $j$. Every step used only
    $j\notin\gamma^{\ast}$ and $\gamma\in\mathscr{M}(s)$: first
    $\Phi_{\gamma}X_{\gamma^{\ast}\cap\gamma}
    =\Phi_{\gamma\setminus\{j\}}X_{\gamma^{\ast}\cap\gamma}$, so the left
    side equals $\|(\Phi_\gamma-\Phi_{\gamma\setminus\{j\}})
    X_{\gamma^{\ast}\setminus\gamma}
    \beta^{\ast}_{\gamma^{\ast}\setminus\gamma}\|_2^{2}$; then the
    identification with a difference of least-squares residuals; then
    feasibility of zeroing the $j$th coordinate, which gives the bound
    $n\abs{\widehat\beta_j(\gamma)}^{2}$; then
    $\|\widehat\beta(\gamma)\|_2^{2}\leq\omega(X)
    \|\beta^{\ast}_{\gamma^{\ast}\setminus\gamma}\|_2^{2}$. Taking
    $\gamma\supseteq\gamma^{\ast}$ makes the right side vanish and
    recovers the identity used in \cref{sub:over}.
\end{proof}

\subsection{Counting and paths}
\label{sub:appcount}

\begin{lemma}[Counting precedents]
    \label{lemma:count}
    Let $\bar\eta\in\mathscr{M}(s_0)$.
    \begin{equivcond}
        \item\label{it:cnt_i} $\mathcal{G}$ maps overfitted states to
        states containing $\gamma^{\ast}$. Consequently, if $\gamma$ is
        underfitted then every $\bar\gamma\in\Lambda(\gamma)$ is
        underfitted.
        \item\label{it:cnt_ii} $\bar\eta$ has at most $(s^{\ast}+1)p$
        immediate precedents that are underfitted --- that is, states
        $\bar\gamma\neq\bar\eta$ with
        $\mathcal{G}(\bar\gamma)=\bar\eta$ and $\bar\gamma$ underfitted
        --- and at most $p$ that are overfitted.
        \item\label{it:cnt_iii} Hence, for every $k\geq0$, the number of
        $\bar\gamma$ with $\mathcal{G}^{k}(\bar\gamma)=\bar\eta$ and
        $\bar\gamma,\mathcal{G}(\bar\gamma),\dots,
        \mathcal{G}^{k-1}(\bar\gamma)$ all underfitted is at most
        $\{(s^{\ast}+1)p\}^{k}$; and if
        $\bar\eta\supseteq\gamma^{\ast}$, the number of overfitted
        $\bar\gamma\in\Lambda_k(\bar\eta)$ is at most $p^{k}$.
    \end{equivcond}
\end{lemma}

\begin{proof}
    \ref{it:cnt_i} An overfitted $\bar\gamma$ is treated by
    rule~\ref{G:del}, which deletes
    $\ell_{\bar\gamma}\in\bar\gamma\setminus\gamma^{\ast}$ and therefore
    leaves $\gamma^{\ast}\subseteq\mathcal{G}(\bar\gamma)$; and
    $\gamma^{\ast}$ is a fixed point. So
    $\mathcal{G}^{k}(\bar\gamma)\supseteq\gamma^{\ast}$ for every
    $k\geq0$, and no such state is underfitted.

    \ref{it:cnt_ii} Let $\bar\gamma\neq\bar\eta$ be underfitted with
    $\mathcal{G}(\bar\gamma)=\bar\eta$. If rule~\ref{G:add} applies then
    $\bar\eta=\bar\gamma\cup\{j_{\bar\gamma}\}$ with
    $j_{\bar\gamma}\in\gamma^{\ast}$, so
    $\bar\gamma=\bar\eta\setminus\{j\}$ for some
    $j\in\gamma^{\ast}\cap\bar\eta$: at most $s^{\ast}$ states. If
    rule~\ref{G:exch} applies then
    $\bar\eta=\bar\gamma\cup\{j_{\bar\gamma}\}
    \setminus\{k_{\bar\gamma}\}$ with $j_{\bar\gamma}\in\gamma^{\ast}$
    and $k_{\bar\gamma}\in S^{c}$, so
    $\bar\gamma=\bar\eta\cup\{k\}\setminus\{j\}$ with
    $j\in\gamma^{\ast}\cap\bar\eta$ and $k\in S^{c}$: at most
    $s^{\ast}(p-s^{\ast})$ states. The total is at most
    $s^{\ast}(1+p-s^{\ast})\leq s^{\ast}(1+p)\leq(s^{\ast}+1)p$, the last
    step because $s^{\ast}\leq p$. If instead $\bar\gamma$ is overfitted
    then rule~\ref{G:del} applies and $\bar\gamma=\bar\eta\cup\{\ell\}$
    with $\ell\in S^{c}$: at most $p$ states.

    \ref{it:cnt_iii} Iterate~\ref{it:cnt_ii}: a state counted at level
    $k$ is an immediate precedent, of the stated type, of one counted at
    level $k-1$. For the second claim, \ref{it:cnt_i} shows that an
    overfitted $\bar\gamma\in\Lambda_k(\bar\eta)$ reaches $\bar\eta$
    through states containing $\gamma^{\ast}$ only, so each of the $k$
    steps deletes an index of $S^{c}$; hence
    $\bar\gamma=\bar\eta\cup u$ with $u\subseteq S^{c}$ and $\abs u=k$,
    and there are at most $\binom{p}{k}\leq p^{k}$ such $u$.
\end{proof}

\begin{lemma}[Path properties; YWJ, Lemma~2]
    \label{lemma:paths}
    \begin{equivcond}
        \item\label{it:pp_i} $\mathcal{G}^{k}(\gamma)=\gamma^{\ast}$ for
        some $k\leq d_H(\gamma,\gamma^{\ast})$, every
        $\gamma\in\mathscr{M}(s_0)$; in particular
        $\Lambda(\gamma^{\ast})=\mathscr{M}(s_0)$.
        \item\label{it:pp_ii} $\abs{T_{\gamma,\gamma^{\ast}}}\leq
        d_H(\gamma,\gamma^{\ast})\leq s_0+s^{\ast}\leq2s_0$, and
        $\ell(\mathcal{T})\leq4s_0$.
        \item\label{it:pp_iii} For a directed edge
        $e=(\bar\gamma,\mathcal{G}(\bar\gamma))$,
        \begin{equation*}
            \bigl\{(\gamma,\gamma'):T_{\gamma,\gamma'}\ni e\bigr\}
            \ \subseteq\ \Lambda(\bar\gamma)\times\mathscr{M}(s_0) .
        \end{equation*}
    \end{equivcond}
\end{lemma}

\begin{proof}
    \ref{it:pp_i} By the second remark after \cref{def:G},
    $d_H(\mathcal{G}(\gamma),\gamma^{\ast})<d_H(\gamma,\gamma^{\ast})$
    whenever $\gamma\neq\gamma^{\ast}$, so the orbit reaches the unique
    fixed point $\gamma^{\ast}$ in at most
    $d_H(\gamma,\gamma^{\ast})$ steps. Hence every state lies in
    $\Lambda(\gamma^{\ast})$.

    \ref{it:pp_ii} Each step of $\mathcal{G}$ decreases
    $d_H(\cdot,\gamma^{\ast})$ by at least one, so the orbit has at most
    $d_H(\gamma,\gamma^{\ast})$ edges; and
    $d_H(\gamma,\gamma^{\ast})\leq\abs\gamma+\abs{\gamma^{\ast}}
    \leq s_0+s^{\ast}$, which is at most $2s_0$ because
    $s^{\ast}<s_0$. Since $T_{\gamma,\gamma'}$ is contained in the
    concatenation of $T_{\gamma,\gamma^{\ast}}$ with the reverse of
    $T_{\gamma',\gamma^{\ast}}$, its length is at most $4s_0$.

    \ref{it:pp_iii} By construction $T_{\gamma,\gamma'}$ traverses its
    first segment in the direction of $\mathcal{G}$ and its second
    segment against it. The edge $e$ points in the direction of
    $\mathcal{G}$, so it can only occur in the first segment, which is
    part of $T_{\gamma,\gamma^{\ast}}$. Hence $\bar\gamma\in
    T_{\gamma,\gamma^{\ast}}$, that is $\gamma\in\Lambda(\bar\gamma)$,
    while $\gamma'$ is unconstrained.
\end{proof}

\subsection{Proof of \cref{th:drop}: overfitted states}
\label{sub:over}

Throughout \cref{sub:over,sub:under,sub:sat} we write
$\gamma'\coloneq\mathcal{G}(\gamma)$ and start from~\eqref{eq:master} with
the size change read off from~\eqref{eq:sizechange}.

Here $\gamma\supseteq\gamma^{\ast}$, $\gamma\neq\gamma^{\ast}$, and
rule~\ref{G:del} applies: $\gamma'=\gamma\setminus\{\ell_\gamma\}$ with
$\ell_\gamma\in\gamma\setminus\gamma^\ast$, so
$\abs{\gamma'}-\abs \gamma=-1$ and \eqref{eq:master} gives
\begin{equation}
    \label{eq:over1}
    \log\frac{\pi_n^{(\lambda)}(\gamma\mid Y)}
    {\pi_n^{(\lambda)}(\gamma'\mid Y)}
    \ =\ -\kappa\log p-\frac{\lambda}{2}\log(1+g)
    +\frac{\lambda n}{2}
    \log\frac{g^{-1}+1-R^{2}_{\gamma'}}{g^{-1}+1-R^{2}_{\gamma}} .
\end{equation}
Since $\gamma'\subset\gamma$ and $\abs \gamma-\abs{\gamma'}=1$, the matrix
$\Phi_\gamma-\Phi_{\gamma'}$ is an orthogonal projection of rank one, and
\begin{equation}
    \label{eq:over1b}
    \norm Y_2^{2}\bigl(R^{2}_{\gamma}-R^{2}_{\gamma'}\bigr)
    =Y^{\top}(\Phi_\gamma-\Phi_{\gamma'})Y
    =\norm{(\Phi_\gamma-\Phi_{\gamma'})Y}_2^{2}\ \geq\ 0 ,
\end{equation}
so that, using $\log(1+x)\leq x$ and then \cref{lemma:residual},
\begin{equation}
    \label{eq:over2}
    \frac n2\log\frac{g^{-1}+1-R^{2}_{\gamma'}}
    {g^{-1}+1-R^{2}_{\gamma}}
    =\frac n2\log\Bigl(1+\frac{\norm{(\Phi_\gamma-\Phi_{\gamma'})
        Y}_2^{2}}{\norm Y_2^{2}(g^{-1}+1-R^{2}_{\gamma})}\Bigr)
    \ \leq\ \frac n2\cdot
    \frac{8\norm{(\Phi_\gamma-\Phi_{\gamma'})Y}_2^{2}}{n\sigma_0^{2}} .
\end{equation}
It remains to bound $\norm{(\Phi_\gamma-\Phi_{\gamma'})Y}_2^{2}$. Since
$\ell_\gamma\notin\gamma^\ast$ and $\gamma^{\ast}\subseteq\gamma$, we have
$\gamma^{\ast}\subseteq\gamma'$, so $\signal$ lies in the range of
$\Phi_{\gamma'}$; and $\Phi_\gamma-\Phi_{\gamma'}$ annihilates that range,
because for $v$ in it $\Phi_\gamma v=v$, the range of $\Phi_{\gamma'}$
being contained in that of $\Phi_\gamma$, so that
$(\Phi_\gamma-\Phi_{\gamma'})v=v-v=0$. Hence, by~\eqref{eq:wtilde},
\begin{equation}
    \label{eq:over3}
    (\Phi_\gamma-\Phi_{\gamma'})Y
    =(\Phi_\gamma-\Phi_{\gamma'})\widetilde w,
    \qquad\text{so}\qquad
    \norm{(\Phi_\gamma-\Phi_{\gamma'})Y}_2^{2}
    \ \leq\ 2(L+\Ltil)\sigma_0^{2}\log p
\end{equation}
by \cref{lemma:noise}\ref{it:noise_ii}, applied to the nested pair
$\gamma'\subset\gamma$ of
$\mathscr{M}(s_0)\subseteq\mathscr{M}(s)$ with
$\abs \gamma-\abs{\gamma'}=1$. Substituting~\eqref{eq:over3}
into~\eqref{eq:over2},
\begin{equation}
    \label{eq:over4}
    \frac n2\log\frac{g^{-1}+1-R^{2}_{\gamma'}}
    {g^{-1}+1-R^{2}_{\gamma}}\ \leq\ 8(L+\Ltil)\log p ,
\end{equation}
and substituting this into~\eqref{eq:over1}, using $\lambda\geq0$ so that
multiplying~\eqref{eq:over4} by $\lambda$ preserves the inequality,
\begin{equation*}
    \log\frac{\pi_n^{(\lambda)}(\gamma\mid Y)}
    {\pi_n^{(\lambda)}(\gamma'\mid Y)}
    \ \leq\ -\kappa\log p
    -\lambda\Bigl\{\frac12\log(1+g)-8(L+\Ltil)\log p\Bigr\} ,
\end{equation*}
which is~\eqref{eq:drop_over}. \qed

\subsection{Proof of \cref{th:drop}: underfitted, unsaturated states}
\label{sub:under}

Here $\gamma^{\ast}\setminus\gamma\neq\emptyset$ and $\abs \gamma<s_0$, so
rule~\ref{G:add} applies: $\gamma'=\gamma\cup\{j_\gamma\}$,
$\abs{\gamma'}-\abs \gamma=+1$, and \eqref{eq:master} gives
\begin{equation}
    \label{eq:und1}
    \log\frac{\pi_n^{(\lambda)}(\gamma\mid Y)}
    {\pi_n^{(\lambda)}(\gamma'\mid Y)}
    \ =\ \kappa\log p+\frac{\lambda}{2}\log(1+g)
    +\frac{\lambda n}{2}
    \log\frac{g^{-1}+1-R^{2}_{\gamma'}}{g^{-1}+1-R^{2}_{\gamma}} .
\end{equation}
Note $s^{\ast}\geq1$, since $\gamma^{\ast}\setminus\gamma\neq\emptyset$.
The matrix $\Phi_{\gamma'}-\Phi_\gamma$ is an orthogonal projection of
rank one, and as in~\eqref{eq:over1b},
\begin{equation}
    \label{eq:und1b}
    \norm Y_2^{2}\bigl(R^{2}_{\gamma'}-R^{2}_{\gamma}\bigr)
    =\norm{(\Phi_{\gamma'}-\Phi_{\gamma})Y}_2^{2} ,
\end{equation}
so that, by $\log(1+x)\leq x$ applied with
$x=-\norm{(\Phi_{\gamma'}-\Phi_\gamma)Y}_2^{2}/
\{\norm Y_2^{2}(g^{-1}+1-R^{2}_\gamma)\}\in(-1,0]$,
\begin{equation}
    \label{eq:und2}
    \frac n2\log\frac{g^{-1}+1-R^{2}_{\gamma'}}
    {g^{-1}+1-R^{2}_{\gamma}}
    \ \leq\ -\frac n2\cdot
    \frac{\norm{(\Phi_{\gamma'}-\Phi_\gamma)Y}_2^{2}}
    {\norm Y_2^{2}\bigl(g^{-1}+1-R^{2}_{\gamma}\bigr)} .
\end{equation}

\emph{Step 1: the numerator.} By \cref{lemma:forward}\ref{it:fs_a},
$\norm{(\Phi_{\gamma'}-\Phi_\gamma)\signal}_2\geq A$, and by
\cref{lemma:noise}\ref{it:noise_ii} applied to the nested pair
$\gamma\subset\gamma'$, together with~\eqref{eq:Alower},
\begin{equation}
    \label{eq:und3}
    \norm{(\Phi_{\gamma'}-\Phi_\gamma)\widetilde w}_2^{2}
    \ \leq\ \frac{n\nu^{2}C_\beta^{2}}{64}\ \leq\ \frac{A^{2}}{64},
    \qquad\text{that is}\qquad
    \norm{(\Phi_{\gamma'}-\Phi_\gamma)\widetilde w}_2
    \leq\frac{A}{8} .
\end{equation}
Hence, by the triangle inequality and~\eqref{eq:wtilde},
\begin{equation}
    \label{eq:und4}
    \norm{(\Phi_{\gamma'}-\Phi_\gamma)Y}_2
    \ \geq\ A-\frac{A}{8}=\frac{7A}{8}\ >\ 0,
    \quad\text{so}\quad
    \norm{(\Phi_{\gamma'}-\Phi_\gamma)Y}_2^{2}
    \ \geq\ \tfrac{49}{64}A^{2}\ \geq\ \tfrac12A^{2} .
\end{equation}

\emph{Step 2: the denominator.} By \cref{lemma:Dupper} and the
definition~\eqref{eq:Adef}, which gives
$\norm{(I-\Phi_\gamma)\signal}_2^{2}
=\abs{\gamma^{\ast}\setminus\gamma}A^{2}/\nu\leq s^{\ast}A^{2}/\nu$,
\begin{equation}
    \label{eq:und5}
    \norm Y_2^{2}\bigl(g^{-1}+1-R^{2}_{\gamma}\bigr)
    \ \leq\ \frac{2s^{\ast}A^{2}}{\nu}+8n\sigma_0^{2} .
\end{equation}

\emph{Step 3: combining.} By~\eqref{eq:und4} and~\eqref{eq:und5}, then by
the third inequality of~\eqref{eq:elementary} with
$a=2s^{\ast}A^{2}/\nu>0$ and $b=8n\sigma_0^{2}>0$,
\begin{equation}
    \label{eq:und6}
    \frac n2\cdot
    \frac{\norm{(\Phi_{\gamma'}-\Phi_\gamma)Y}_2^{2}}
    {\norm Y_2^{2}(g^{-1}+1-R^{2}_{\gamma})}
    \ \geq\ \frac n4\cdot\frac{A^{2}}
    {2s^{\ast}A^{2}/\nu+8n\sigma_0^{2}}
    \ =\ \frac{n\nu}{8s^{\ast}}\cdot\frac{a}{a+b}
    \ \geq\ \frac{n\nu}{8s^{\ast}}
    \min\Bigl\{\frac12,\frac{s^{\ast}A^{2}}
    {8\nu n\sigma_0^{2}}\Bigr\} ,
\end{equation}
that is, distributing $n\nu/(8s^{\ast})$ over the minimum and then
using $A^{2}\geq n\nu^{2}C_\beta^{2}$ from~\eqref{eq:Alower},
\begin{equation}
    \label{eq:und8}
    \frac n2\cdot
    \frac{\norm{(\Phi_{\gamma'}-\Phi_\gamma)Y}_2^{2}}
    {\norm Y_2^{2}(g^{-1}+1-R^{2}_{\gamma})}
    \ \geq\ \min\Bigl\{\frac{n\nu}{16s^{\ast}},\
    \frac{A^{2}}{64\,\sigma_0^{2}}\Bigr\}
    \ \geq\ \gainfloor .
\end{equation}
Combining \eqref{eq:und1}, \eqref{eq:und2} and~\eqref{eq:und8}, and using
$\lambda\geq0$,
\begin{equation*}
    \log\frac{\pi_n^{(\lambda)}(\gamma\mid Y)}
    {\pi_n^{(\lambda)}(\gamma'\mid Y)}
    \ \leq\ \kappa\log p
    +\lambda\Bigl\{\frac12\log(1+g)-\gainfloor\Bigr\} ,
\end{equation*}
which is~\eqref{eq:drop_unsat}. \qed

\subsection{Proof of \cref{th:drop}: underfitted, saturated states}
\label{sub:sat}

Here $\gamma^{\ast}\setminus\gamma\neq\emptyset$ and $\abs \gamma=s_0$, so
rule~\ref{G:exch} applies: with
$\widetilde\gamma=\gamma\cup\{j_\gamma\}$ and
$\gamma'=\widetilde\gamma\setminus\{k_\gamma\}$ we have
$\abs{\gamma'}=\abs \gamma$, and \eqref{eq:master} gives
\begin{equation}
    \label{eq:sat1}
    \log\frac{\pi_n^{(\lambda)}(\gamma\mid Y)}
    {\pi_n^{(\lambda)}(\gamma'\mid Y)}
    \ =\ \frac{\lambda n}{2}
    \log\frac{g^{-1}+1-R^{2}_{\gamma'}}{g^{-1}+1-R^{2}_{\gamma}} ,
\end{equation}
the inclusion penalty having cancelled identically. Note that
$\widetilde\gamma\in\mathscr{M}(s)$, since
$\abs{\widetilde\gamma}=s_0+1\leq s$, and that both $\gamma$ and $\gamma'$
are subsets of $\widetilde\gamma$ of cardinality
$\abs{\widetilde\gamma}-1$, so that
$\Phi_{\widetilde\gamma}-\Phi_\gamma$ and
$\Phi_{\widetilde\gamma}-\Phi_{\gamma'}$ are orthogonal projections of
rank one. Following YWJ, set
\begin{equation}
    \label{eq:vdef}
    v_1\coloneq(\Phi_{\widetilde\gamma}-\Phi_\gamma)\signal,
    \qquad
    v_2\coloneq(\Phi_{\widetilde\gamma}-\Phi_{\gamma'})\signal .
\end{equation}
Adding and subtracting $Y^{\top}\Phi_{\widetilde\gamma}Y$,
\begin{equation}
    \label{eq:sat2}
    \norm Y_2^{2}\bigl(R^{2}_{\gamma'}-R^{2}_{\gamma}\bigr)
    =Y^{\top}\Phi_{\gamma'}Y-Y^{\top}\Phi_{\gamma}Y
    =\norm{(\Phi_{\widetilde\gamma}-\Phi_\gamma)Y}_2^{2}
    -\norm{(\Phi_{\widetilde\gamma}-\Phi_{\gamma'})Y}_2^{2} .
\end{equation}
By \cref{lemma:forward}\ref{it:fs_a} and~\ref{it:fs_b},
$\norm{v_1}_2\geq A$ and $\norm{v_2}_2\leq A/2$; and by
\cref{lemma:noise}\ref{it:noise_ii} applied to each of the two nested
pairs $\gamma\subset\widetilde\gamma$ and
$\gamma'\subset\widetilde\gamma$ of $\mathscr{M}(s)$, exactly as
in~\eqref{eq:und3},
\begin{equation}
    \label{eq:sat4}
    \norm{(\Phi_{\widetilde\gamma}-\Phi_\gamma)\widetilde w}_2
    \ \leq\ \frac{A}{8},
    \qquad
    \norm{(\Phi_{\widetilde\gamma}-\Phi_{\gamma'})\widetilde w}_2
    \ \leq\ \frac{A}{8} .
\end{equation}
Therefore, by the triangle inequality applied to
$(\Phi_{\widetilde\gamma}-\Phi_\gamma)Y=v_1
+(\Phi_{\widetilde\gamma}-\Phi_\gamma)\widetilde w$ and to
$(\Phi_{\widetilde\gamma}-\Phi_{\gamma'})Y=v_2
+(\Phi_{\widetilde\gamma}-\Phi_{\gamma'})\widetilde w$,
\begin{equation}
    \label{eq:sat5}
    \norm{(\Phi_{\widetilde\gamma}-\Phi_\gamma)Y}_2
    \ \geq\ A-\frac{A}{8}=\frac{7A}{8}\ >\ 0,
    \qquad
    \norm{(\Phi_{\widetilde\gamma}-\Phi_{\gamma'})Y}_2
    \ \leq\ \frac{A}{2}+\frac{A}{8}=\frac{5A}{8} ,
\end{equation}
and squaring, which is legitimate because both bounds are non-negative,
\begin{equation}
    \label{eq:satdrop}
    \norm Y_2^{2}\bigl(R^{2}_{\gamma'}-R^{2}_{\gamma}\bigr)
    \ \geq\ \frac{49}{64}A^{2}-\frac{25}{64}A^{2}
    \ =\ \frac{24}{64}A^{2}\ \geq\ \frac{A^{2}}{4}\ >\ 0 .
\end{equation}
Inequality~\eqref{eq:und2} holds verbatim, with
$\norm{(\Phi_{\gamma'}-\Phi_\gamma)Y}_2^{2}$ replaced by the left-hand
side of~\eqref{eq:satdrop}, and repeating Steps~2 and~3 of
\cref{sub:under} with $A^{2}/2$ replaced by $A^{2}/4$ in the numerator
of~\eqref{eq:und6} halves the right-hand side of~\eqref{eq:und8}:
\begin{equation*}
    \frac n2\cdot
    \frac{\norm Y_2^{2}(R^{2}_{\gamma'}-R^{2}_{\gamma})}
    {\norm Y_2^{2}(g^{-1}+1-R^{2}_{\gamma})}
    \ \geq\ \frac n8\cdot\frac{A^{2}}
    {2s^{\ast}A^{2}/\nu+8n\sigma_0^{2}}
    \ \geq\ \frac12\,\gainfloor .
\end{equation*}
Substituting into~\eqref{eq:sat1} and using $\lambda\geq0$ gives the third
bound~\eqref{eq:drop_sat}. \qed

\subsection{Proofs for \cref{sec:ywj}}
\label{sub:appywj}

\begin{proof}[Proof of \cref{lemma:kernel}]
    Write $q(\gamma,\gamma')$ for the probability that
    $\widetilde{\mathbf{P}}$ proposes $\gamma'$ from $\gamma$. On
    $\mathcal{N}_1$ it equals $1/(2p)$, symmetric in $(\gamma,\gamma')$.
    On $\mathcal{N}_2$ the move preserves the size, so
    $\abs{\gamma'}=\abs\gamma$ and $q(\gamma,\gamma')
    =\{2\abs\gamma(p-\abs\gamma)\}^{-1}=q(\gamma',\gamma)$; in particular
    an exchange never leaves $\mathscr{M}(s_0)$. Hence $q$ is symmetric,
    and for $\gamma\neq\gamma'$
    \begin{equation}
        \label{eq:detbal}
        \mu(\gamma)\,\widetilde{\mathbf{P}}(\gamma,\gamma')
        =q(\gamma,\gamma')\min\bigl\{\mu(\gamma),\mu(\gamma')\bigr\},
    \end{equation}
    which is symmetric; this is detailed balance for
    $\widetilde{\mathbf{P}}$, hence for $\mathbf{P}$, and
    it shows $\mathbf{Q}$ is symmetric.

    For irreducibility, recall that $\mu$ has full support on
    $\mathscr{M}(s_0)$. Given
    $\gamma\neq\gamma'$, delete the elements of $\gamma$ one at a time
    and then add those of $\gamma'$ one at a time; every intermediate
    state is a subset of $\gamma$ or of $\gamma'$, hence lies in
    $\mathscr{M}(s_0)$, and consecutive states differ by one flip, so
    each transition has positive probability by~\eqref{eq:MH}.
    Aperiodicity is $\mathbf{P}(\gamma,\gamma)\geq1/2>0$.

    For non-negative definiteness, $\widetilde{\mathbf{P}}$ is a
    reversible Markov kernel, hence self-adjoint on
    $\mathcal{L}^2(\mu)$ with spectrum in $[-1,1]$, so
    $\mathbf{P}=\tfrac12(\mathbf{I}+\widetilde{\mathbf{P}})$ has spectrum
    in $[0,1]$ and its absolute spectral gap coincides with $1-\lambda_2$.

    For~\eqref{eq:edge}, combine~\eqref{eq:lazy} and~\eqref{eq:detbal}:
    $\mathbf{Q}(\gamma,\gamma')=\tfrac12 q(\gamma,\gamma')
    \min\{\mu(\gamma),\mu(\gamma')\}$, and
    $q(\gamma,\gamma')\geq1/(2ps_0)$ in both cases: on $\mathcal{N}_1$
    because $1/(2p)\geq1/(2ps_0)$, and on $\mathcal{N}_2$ because
    $\abs\gamma\leq s_0$ and $p-\abs\gamma\leq p$.
\end{proof}

\subsection{Proofs for \cref{sec:main}}
\label{sub:appmain}

\begin{proof}[Proof of \cref{cor:critical}]
    For an underfitted unsaturated $\gamma$, the bound
    \eqref{eq:drop_unsat} is at most $-\eta\log p$ precisely when
    $\lambda\{\gainfloorsm-\tfrac12\log(1+g)\}\geq(\kappa+\eta)\log p$,
    that is when $\lambda$ is at least the first term
    of~\eqref{eq:lambdastar}, the coefficient
    $\gainfloorsm-\tfrac12\log(1+g)$ being positive by hypothesis. For an
    underfitted saturated $\gamma$, the bound~\eqref{eq:drop_sat} is
    at most $-\eta\log p$ precisely when
    $\lambda\gainfloorsm\geq2\eta\log p$, that is when $\lambda$ is at
    least the second term of~\eqref{eq:lambdastar}. Taking the maximum
    gives~\eqref{eq:underconc}. The overfitted claim restates
    \eqref{eq:drop_over}, and~\eqref{eq:endpoints} is
    obtained by evaluating the affine function at $\lambda=0$ and
    $\lambda=1$ and bounding $\tfrac12\log(1+g)$ below by
    $\alpha\log p$.
\end{proof}

\begin{proof}[Proof of \cref{cor:prec}]
    Write $a\coloneq(s^{\ast}+1)p^{1-\eta}$ and $b\coloneq p^{1-\eta}$,
    so that $b\leq a\leq\tfrac12$ by~\eqref{eq:etamin}.

    \emph{Step 0: the drop along a path.} If
    $\mathcal{G}^{k}(\bar\gamma)=\gamma$ then, telescoping and applying
    \cref{cor:critical} at each of the $k$ source states --- each of
    which is either overfitted or underfitted, and both cases are covered
    by its hypotheses ---
    \begin{equation}
        \label{eq:telescope}
        \frac{\pi_n^{(\lambda)}(\bar\gamma\mid Y)}
        {\pi_n^{(\lambda)}(\gamma\mid Y)}
        =\prod_{i=0}^{k-1}
        \frac{\pi_n^{(\lambda)}(\mathcal{G}^{i}(\bar\gamma)\mid Y)}
        {\pi_n^{(\lambda)}(\mathcal{G}^{i+1}(\bar\gamma)\mid Y)}
        \ \leq\ p^{-\eta k} .
    \end{equation}

    \emph{Step 1: $\gamma$ underfitted.} By
    \cref{lemma:count}\ref{it:cnt_i} every $\bar\gamma\in\Lambda(\gamma)$
    is underfitted, so \cref{lemma:count}\ref{it:cnt_iii} gives
    $\abs{\Lambda_k(\gamma)}\leq\{(s^{\ast}+1)p\}^{k}$, and
    with~\eqref{eq:telescope}
    \begin{equation}
        \label{eq:precunder}
        \frac{\pi_n^{(\lambda)}[\Lambda(\gamma)\mid Y]}
        {\pi_n^{(\lambda)}(\gamma\mid Y)}
        \ \leq\ \sum_{k\geq0}\{(s^{\ast}+1)p\}^{k}p^{-\eta k}
        \ =\ \sum_{k\geq0}a^{k}\ =\ \frac{1}{1-a} ,
    \end{equation}
    the series converging because $a\leq\tfrac12$.

    \emph{Step 2: $\gamma\supseteq\gamma^{\ast}$.} Split
    $\Lambda(\gamma)=\mathbb{M}_1\sqcup\mathbb{M}_2$ into the precedents
    that contain $\gamma^{\ast}$ and those that do not, so that
    $\gamma\in\mathbb{M}_1$ and every member of $\mathbb{M}_2$ is
    underfitted. By \cref{lemma:count}\ref{it:cnt_iii}
    and~\eqref{eq:telescope},
    \begin{equation}
        \label{eq:precM1}
        \frac{\pi_n^{(\lambda)}[\mathbb{M}_1\mid Y]}
        {\pi_n^{(\lambda)}(\gamma\mid Y)}
        \ \leq\ \sum_{k\geq0}p^{k}p^{-\eta k}\ =\ \frac{1}{1-b} .
    \end{equation}
    For $\bar\gamma\in\mathbb{M}_2$ let $f(\bar\gamma)$ be the first
    state on the path
    $\bar\gamma\to\mathcal{G}(\bar\gamma)\to\dots\to\gamma$ that contains
    $\gamma^{\ast}$; it exists because $\gamma$ itself does, and it lies
    in $\mathbb{M}_1$ because it is a precedent of $\gamma$. Every state
    strictly before $f(\bar\gamma)$ on that path is underfitted, by the
    definition of $f$, so \cref{lemma:count}\ref{it:cnt_iii} applies to
    the chain from $\bar\gamma$ to $f(\bar\gamma)$, whose length is at
    least $1$ because $\bar\gamma\notin\mathbb{M}_1$. Grouping the sum
    over $\mathbb{M}_2$ by the value of $f$, and
    using~\eqref{eq:telescope} for the inner ratios,
    \begin{align}
        \notag
        \frac{\pi_n^{(\lambda)}[\mathbb{M}_2\mid Y]}
        {\pi_n^{(\lambda)}(\gamma\mid Y)}
        &=\sum_{\widetilde\gamma\in\mathbb{M}_1}
        \frac{\pi_n^{(\lambda)}(\widetilde\gamma\mid Y)}
        {\pi_n^{(\lambda)}(\gamma\mid Y)}
        \sum_{\bar\gamma:\,f(\bar\gamma)=\widetilde\gamma}
        \frac{\pi_n^{(\lambda)}(\bar\gamma\mid Y)}
        {\pi_n^{(\lambda)}(\widetilde\gamma\mid Y)}\\
        \label{eq:precM2}
        &\ \leq\ \frac{\pi_n^{(\lambda)}[\mathbb{M}_1\mid Y]}
        {\pi_n^{(\lambda)}(\gamma\mid Y)}
        \cdot\sum_{k\geq1}a^{k}
        \ \leq\ \frac{1}{1-b}\cdot\frac{a}{1-a} .
    \end{align}
    Adding~\eqref{eq:precM1} and~\eqref{eq:precM2},
    \begin{equation*}
        \frac{\pi_n^{(\lambda)}[\Lambda(\gamma)\mid Y]}
        {\pi_n^{(\lambda)}(\gamma\mid Y)}
        \ \leq\ \frac{1}{1-b}\Bigl(1+\frac{a}{1-a}\Bigr)
        \ =\ \frac{1}{(1-a)(1-b)} ,
    \end{equation*}
    which is the middle expression in~\eqref{eq:precmass}. It dominates
    the bound $1/(1-a)$ of Step~1, so both cases are covered by it, and
    every $\gamma\in\mathscr{M}(s_0)$ falls into one of them: a state
    other than $\gamma^{\ast}$ is either underfitted or overfitted, and
    $\gamma^{\ast}\supseteq\gamma^{\ast}$. Finally $a,b\leq\tfrac12$
    gives $\{(1-a)(1-b)\}^{-1}\leq4$.
\end{proof}

\begin{proof}[Proof of \cref{prop:sharp}]
    Identity~\eqref{eq:sharpid} is~\eqref{eq:onestep} with
    $j=j_\gamma$, read in the direction
    $\log\{\pi_n^{(\lambda)}(\gamma\mid Y)/
    \pi_n^{(\lambda)}(\gamma'\mid Y)\}
    =-\log\{\pi_n^{(\lambda)}(\gamma'\mid Y)/
    \pi_n^{(\lambda)}(\gamma\mid Y)\}$. The right-hand side
    of~\eqref{eq:sharpid} is non-positive if and only if
    $\lambda\log\{\mathcal{L}_n(Y\mid\mathcal{G}(\gamma))/
    \mathcal{L}_n(Y\mid\gamma)\}\geq\kappa\log p$, which gives the stated
    equivalence; at $\lambda=0$ it equals $\kappa\log p>0$; and if the
    likelihood ratio is at most one then its logarithm is non-positive, so
    for $\lambda\geq0$ the right-hand side is at least $\kappa\log p>0$.
    The characterisation of $\lambda_{\mathrm{crit}}$ follows by taking
    the maximum over the finitely many underfitted unsaturated $\gamma$,
    with the convention that a state whose likelihood ratio is at most one
    contributes $+\infty$.

    For~\eqref{eq:lambdacritbound}, form the ratio of two copies
    of~\eqref{eq:Lik} with $\gamma'=\gamma\cup\{j_\gamma\}$:
    \begin{align*}
        \log\frac{\mathcal{L}_n(Y\mid\gamma')}
        {\mathcal{L}_n(Y\mid\gamma)}
        &=-\frac12\log(1+g)
        -\frac n2\log\frac{g^{-1}+1-R^{2}_{\gamma'}}
        {g^{-1}+1-R^{2}_{\gamma}}\\
        &\ \geq\ -\frac12\log(1+g)+\gainfloor ,
    \end{align*}
    the last step by~\eqref{eq:und2} and~\eqref{eq:und8}. The right-hand
    side being positive by hypothesis, the quotient defining
    $\lambda_{\mathrm{crit}}$ is at most
    $\kappa\log p/\{\gainfloorsm-\tfrac12\log(1+g)\}$ for every such
    $\gamma$, and taking the maximum gives the claim.
\end{proof}

\subsection{Proofs for \cref{sec:gap}}
\label{sub:appgap}

\begin{proof}[Proof of \cref{lemma:gap}]
    By \cref{cor:critical} the drop bound gives
    $\pi_n^{(\lambda)}(\bar\gamma\mid Y)\leq
    \pi_n^{(\lambda)}(\mathcal{G}(\bar\gamma)\mid Y)$ at every
    $\bar\gamma\neq\gamma^{\ast}$ --- each such state being either
    overfitted or underfitted, and both cases being covered by its
    hypotheses --- so the inner maximum in~\eqref{eq:rhoreduce} equals
    one, and \cref{cor:prec} bounds the remaining factor by $4$. Hence
    $\rho(\mathcal{T})\leq4ps_0\cdot4=16ps_0$. With
    $\ell(\mathcal{T})\leq4s_0$ from
    \cref{lemma:paths}\ref{it:pp_ii}, Sinclair's
    bound~\eqref{eq:sinclair} gives
    $\mathrm{Gap}(\mathbf{P}_\lambda)\geq(16ps_0\cdot4s_0)^{-1}$, which
    is~\eqref{eq:gap}; it is the absolute spectral gap by
    \cref{lemma:kernel}.
\end{proof}

\begin{proof}[Proof of \cref{prop:conc}]
    \emph{Step 1: the precedents of $\gamma^{\ast}$ are everything.} By
    \cref{lemma:paths}\ref{it:pp_i}, every $\bar\gamma\in\mathscr{M}(s_0)$
    satisfies $\mathcal{G}^{k}(\bar\gamma)=\gamma^{\ast}$ for some
    $k\geq0$, which is the defining condition~\eqref{eq:precdef} for
    membership of $\Lambda(\gamma^{\ast})$. Hence
    $\Lambda(\gamma^{\ast})\supseteq\mathscr{M}(s_0)$, and the reverse
    inclusion holds by~\eqref{eq:precdef}, so
    $\Lambda(\gamma^{\ast})=\mathscr{M}(s_0)$. Since
    $\pi_n^{(\lambda)}(\cdot\mid Y)$ is a probability measure on
    $\mathscr{M}(s_0)$ by~\eqref{eq:pilambda},
    \begin{equation}
        \label{eq:concmass}
        \pi_n^{(\lambda)}\bigl[\Lambda(\gamma^{\ast})\mid Y\bigr]
        \ =\ \pi_n^{(\lambda)}\bigl[\mathscr{M}(s_0)\mid Y\bigr]
        \ =\ 1 .
    \end{equation}

    \emph{Step 2: \cref{cor:prec} applies at $\gamma=\gamma^{\ast}$.}
    That corollary is stated for every $\gamma\in\mathscr{M}(s_0)$, and
    $\gamma^{\ast}$ is covered by Step~2 of its proof, which treats every
    $\gamma$ with $\gamma\supseteq\gamma^{\ast}$; the case
    $\gamma=\gamma^{\ast}$ is admissible there because
    $\gamma^{\ast}\supseteq\gamma^{\ast}$. Substituting
    \eqref{eq:concmass} into~\eqref{eq:precmass} therefore gives
    \begin{equation}
        \label{eq:concinv}
        \frac{1}{\pi_n^{(\lambda)}(\gamma^{\ast}\mid Y)}
        \ =\ \frac{\pi_n^{(\lambda)}[\Lambda(\gamma^{\ast})\mid Y]}
        {\pi_n^{(\lambda)}(\gamma^{\ast}\mid Y)}
        \ \leq\ \frac{1}
        {\bigl\{1-(s^{\ast}+1)p^{1-\eta}\bigr\}
        \bigl\{1-p^{1-\eta}\bigr\}} .
    \end{equation}

    \emph{Step 3: inverting.} Put
    $u\coloneq\bigl\{1-(s^{\ast}+1)p^{1-\eta}\bigr\}
    \bigl\{1-p^{1-\eta}\bigr\}$, the right-hand side
    of~\eqref{eq:conc}. By~\eqref{eq:etamin} we have
    $p^{1-\eta}\leq(s^{\ast}+1)p^{1-\eta}\leq\tfrac12$, so both factors
    of $u$ lie in $[\tfrac12,1]$ and hence
    \begin{equation}
        \label{eq:urange}
        \tfrac14\ \leq\ u\ \leq\ 1 .
    \end{equation}
    Also $\pi_n^{(\lambda)}(\gamma^{\ast}\mid Y)\in(0,1]$
    by~\eqref{eq:pilambda}. Inequality~\eqref{eq:concinv} therefore reads
    \begin{equation*}
        \frac{1}{\pi_n^{(\lambda)}(\gamma^{\ast}\mid Y)}
        \ \leq\ \frac{1}{u},
    \end{equation*}
    an inequality between two elements of $(0,\infty)$. The map
    $x\mapsto1/x$ is strictly decreasing on $(0,\infty)$, so applying it
    to both sides reverses the inequality and gives
    $\pi_n^{(\lambda)}(\gamma^{\ast}\mid Y)\geq u$, which is the first
    inequality of~\eqref{eq:conc}. The second is the left-hand
    inequality of~\eqref{eq:urange}.
\end{proof}

\begin{proof}[Proof of \cref{lemma:minmass}]
    Apply \eqref{eq:master} with $\gamma'=\gamma^{\ast}$:
    \begin{equation*}
        \log\frac{\pi_n^{(\lambda)}(\gamma\mid Y)}
        {\pi_n^{(\lambda)}(\gamma^{\ast}\mid Y)}
        =\bigl(\abs{\gamma^{\ast}}-\abs{\gamma}\bigr)
        \Bigl\{\kappa\log p+\frac{\lambda}{2}\log(1+g)\Bigr\}
        +\frac{\lambda n}{2}
        \log\frac{g^{-1}+1-R^{2}_{\gamma^{\ast}}}
        {g^{-1}+1-R^{2}_{\gamma}} .
    \end{equation*}
    For the first term, $\abs{\gamma^{\ast}}-\abs\gamma\geq-\abs\gamma
    \geq-s_0$ and the brace is non-negative, so it is at least
    $-s_0\{\kappa\log p+\tfrac{\lambda}{2}\log(1+g)\}$. For the second,
    $R^{2}_{\gamma^{\ast}}\leq1$ and $R^{2}_{\gamma}\geq0$ give
    \begin{equation*}
        \frac{g^{-1}+1-R^{2}_{\gamma^{\ast}}}
        {g^{-1}+1-R^{2}_{\gamma}}
        \ \geq\ \frac{g^{-1}}{g^{-1}+1}\ =\ \frac{1}{1+g},
    \end{equation*}
    so it is at least $-\tfrac{\lambda n}{2}\log(1+g)$. Adding the two
    and using $\lambda\geq0$ gives~\eqref{eq:minmass}. For
    \eqref{eq:minmass2}, combine it with $\pi_n^{(\lambda)}
    (\gamma^{\ast}\mid Y)\geq\tfrac14$ from \cref{prop:conc}.
\end{proof}

\begin{proof}[Proof of \cref{th:mixing}]
    \cref{lemma:kernel} makes $\mathbf{P}_\lambda$ lazy and reversible,
    so~\eqref{eq:sandwich} applies with
    $\mu=\pi_n^{(\lambda)}(\cdot\mid Y)$. Substitute the gap bound
    of \cref{lemma:gap} and the mass bound~\eqref{eq:minmass2}, whose
    logarithm is $\kappa s_0\log p+\tfrac{\lambda(n+s_0)}{2}\log(1+g)
    +\log4$.
\end{proof}

\subsection{Proofs for \cref{sec:alltemp}}
\label{sub:appall}

\begin{proof}[Proof of \cref{lemma:Gvalid}]
    \ref{it:gv_i} Rules \ref{H:null}, \ref{H:del} and \ref{H:undec}
    delete a coordinate, so they decrease $\abs\gamma$ and stay in
    $\mathscr{M}(s_0)$. Rule~\ref{H:add} fires only when
    $\gamma\subseteq\gamma^{\ast}$, so the result is a subset of
    $\gamma^{\ast}$ and has size at most $s^{\ast}<s_0$.

    \ref{it:gv_ii} Set
    \begin{equation}
        \label{eq:Vprog}
        V(\gamma)\coloneq\abs{\gamma\setminus\gamma^{\ast}}
        +\abs{c^{-}_\lambda\setminus\gamma}
        +\abs{(\gamma\cap\gamma^{\ast})\setminus c^{+}_\lambda}
        +\abs{\gamma\cap C_\lambda} .
    \end{equation}
    Rule~\ref{H:null} decreases the first term by one and leaves the
    others unchanged, since it deletes an index of $S^{c}$ while
    $c^{-}_\lambda,c^{+}_\lambda,C_\lambda\subseteq\gamma^{\ast}$. Rule
    \ref{H:add} decreases the second by one and leaves the others
    unchanged, since $c^{-}_\lambda\cap C_\lambda=\emptyset$ and
    $c^{-}_\lambda\subseteq c^{+}_\lambda$. Rule~\ref{H:del} decreases
    the third by one and leaves the others unchanged, the deleted index
    lying outside $c^{+}_\lambda\supseteq C_\lambda\cup c^{-}_\lambda$.
    Rule~\ref{H:undec} decreases the fourth by one and leaves the others
    unchanged, the deleted index lying in
    $C_\lambda\subseteq c^{+}_\lambda\setminus c^{-}_\lambda$. So $V$
    drops by exactly one at each step, and $V(\gamma)=0$ forces
    $\gamma\subseteq\gamma^{\ast}$, $c^{-}_\lambda\subseteq\gamma$,
    $\gamma\subseteq c^{+}_\lambda$ and $\gamma\cap C_\lambda=\emptyset$,
    that is $\gamma=c^{-}_\lambda$; conversely
    $V(c^{-}_\lambda)=0$ and rule~\ref{H:fix} fixes it. Finally
    $V(\gamma)\leq\abs{\gamma}+\abs{c^{-}_\lambda}\leq s_0+s^{\ast}
    \leq2s_0$, because the first, third and fourth terms
    of~\eqref{eq:Vprog} count disjoint subsets of $\gamma$.

    \ref{it:gv_iii} Rules \ref{H:add}, \ref{H:del} and \ref{H:undec} fire
    only when $\gamma\subseteq\gamma^{\ast}$ and they keep the state
    inside $\gamma^{\ast}$, so \ref{H:null} cannot fire afterwards.
    Rule~\ref{H:undec} deletes an index of $C_\lambda$, which changes
    neither $c^{-}_\lambda\subseteq\gamma$ nor
    $\gamma\subseteq c^{+}_\lambda$, so once those two conditions hold
    they continue to hold.
\end{proof}

\begin{proof}[Proof of \cref{th:dropG}]
    All three are~\eqref{eq:onestepG}.

    \emph{Rule~\ref{H:null}.} Here $j=\max(\gamma\setminus\gamma^{\ast})
    \in S^{c}$ and $\mathcal{G}_\lambda(\gamma)=\gamma\setminus\{j\}$, so
    \eqref{eq:onestepG} applied at the state $\gamma\setminus\{j\}$,
    whose size is $<s_0$, gives
    \begin{equation*}
        \log\frac{\pi_n^{(\lambda)}(\gamma\mid Y)}
        {\pi_n^{(\lambda)}(\gamma\setminus\{j\}\mid Y)}
        =\lambda\psi_j(\gamma\setminus\{j\})-\kappa\log p
        \ \leq\ \lambda(a_0-\alpha)\log p-\kappa\log p
        =-\kappa_\lambda\log p
    \end{equation*}
    by \cref{ass:null} and $\lambda\geq0$; and
    $\kappa_\lambda\geq\eta$ by~\eqref{eq:kappalam}.

    \emph{Rules~\ref{H:add} and~\ref{H:del}.} Here
    $\gamma\subseteq\gamma^{\ast}$, and the state at which
    \eqref{eq:onestepG} is applied --- $\gamma$ under~\ref{H:add},
    $\gamma\setminus\{j\}$ under~\ref{H:del} --- is contained in
    $\gamma^{\ast}\setminus\{j\}$. Hence \eqref{eq:cpmuse} applies, with
    $j\in c^{-}_\lambda$ in the first case and $j\notin c^{+}_\lambda$ in
    the second, and in both
    $\log\{\pi_n^{(\lambda)}(\gamma\mid Y)/
    \pi_n^{(\lambda)}(\mathcal{G}_\lambda(\gamma)\mid Y)\}
    \leq-\log(2s^{\ast})$.

    \emph{Rule~\ref{H:undec}.} Identical to rule~\ref{H:del}, with
    $j\in C_\lambda$ and without the conclusion
    $\lambda\psi_j^{+}\leq\kappa\log p-\log(2s^{\ast})$, which is
    precisely what fails for an undecided covariate; the bound
    $\exp\{\lambda\psi_j^{+}-\kappa\log p\}\leq\Theta_\lambda$ is all that
    remains.
\end{proof}

\begin{proof}[Proof of \cref{lemma:countG}]
    Under~\ref{H:null}, $\bar\gamma=\bar\eta\cup\{\ell\}$ with
    $\ell\in S^{c}$: at most $p$ choices. Under~\ref{H:add},
    $\bar\gamma=\bar\eta\setminus\{j\}$ with
    $j\in c^{-}_\lambda\subseteq\gamma^{\ast}$: at most $s^{\ast}$.
    Under~\ref{H:del}, $\bar\gamma=\bar\eta\cup\{j\}$ with
    $j\in\gamma^{\ast}\setminus c^{+}_\lambda$: at most $s^{\ast}$.
    Under~\ref{H:undec}, $\bar\gamma=\bar\eta\cup\{j\}$ with
    $j\in C_\lambda$: at most $k_\lambda$.
\end{proof}

\begin{proof}[Proof of \cref{cor:precG}]
    By \cref{lemma:Gvalid}\ref{it:gv_iii} the orbit of any $\bar\gamma$
    consists of a block of \ref{H:null} steps, then a block of
    \ref{H:add}/\ref{H:del} steps, then a block of \ref{H:undec} steps.
    Reading this backwards from $\gamma$, a precedent
    $\bar\gamma\in\Lambda_\lambda(\gamma)$ is obtained from $\gamma$ by
    undoing $m$ steps of type \ref{H:undec}, then $k_2$ of type
    \ref{H:add}/\ref{H:del}, then $k_1$ of type \ref{H:null}, for a
    unique triple $(m,k_2,k_1)$, and by \cref{lemma:Gvalid}\ref{it:gv_ii}
    the deletions of \ref{H:undec} and the flips of \ref{H:add},
    \ref{H:del} exhaust $C_\lambda$ and
    $\gamma^{\ast}\triangle c^{-}_\lambda$ respectively, so
    $m\leq k_\lambda$ and $k_2\leq2s^{\ast}$.

    Telescoping as in~\eqref{eq:telescope} and applying \cref{th:dropG}
    at each source state, the ratio
    $\pi_n^{(\lambda)}(\bar\gamma\mid Y)/\pi_n^{(\lambda)}(\gamma\mid Y)$
    is at most $\Theta_\lambda^{m}(2s^{\ast})^{-k_2}p^{-\eta k_1}$, while
    by \cref{lemma:countG} the number of such $\bar\gamma$ is at most
    $k_\lambda^{m}(2s^{\ast})^{k_2}p^{k_1}$. Summing over the three
    indices independently,
    \begin{align*}
        \frac{\pi_n^{(\lambda)}[\Lambda_\lambda(\gamma)\mid Y]}
        {\pi_n^{(\lambda)}(\gamma\mid Y)}
        &\ \leq\
        \Bigl(\sum_{m=0}^{k_\lambda}(k_\lambda\Theta_\lambda)^{m}\Bigr)
        \Bigl(\sum_{k_2=0}^{2s^{\ast}}1\Bigr)
        \Bigl(\sum_{k_1\geq0}p^{(1-\eta)k_1}\Bigr)\\
        &\ \leq\ \bigl(1+k_\lambda\Theta_\lambda\bigr)^{k_\lambda}
        \,(2s^{\ast}+1)\,\frac{1}{1-p^{1-\eta}},
    \end{align*}
    where the first sum was bounded by the binomial theorem, the second
    counts $k_2\leq2s^{\ast}$ terms each equal to one, and the third
    converges because $p^{1-\eta}\leq\tfrac12$, which also bounds it by
    $2$.
\end{proof}

\begin{proof}[Proof of Proposition~\ref{prop:a0}]
    Fix $\gamma$ and $j\in S^{c}\setminus\gamma$ and write
    $\Pi=\Phi_{\gamma\cup\{j\}}-\Phi_\gamma$. By~\eqref{eq:Lik},
    $\psi_j(\gamma)=-\tfrac12\log(1+g)+\tfrac n2\log\{(g^{-1}
    +1-R^{2}_{\gamma})/(g^{-1}+1-R^{2}_{\gamma\cup\{j\}})\}$, and the
    computation~\eqref{eq:over1b}--\eqref{eq:over2}, which used only that
    $\Pi$ is a rank-one projection and \cref{lemma:residual}, gives
    \begin{equation*}
        \psi_j(\gamma)\ \leq\ -\alpha\log p
        +\frac{4\norm{\Pi Y}_2^{2}}{\sigma_0^{2}} ,
    \end{equation*}
    using $\tfrac12\log(1+g)\geq\alpha\log p$ from \cref{ass:C}. Now
    $\Pi Y=\Pi\signal+\Pi\widetilde w$ by~\eqref{eq:wtilde}, so
    by~\eqref{eq:elementary}, \cref{lemma:proxy} applied to
    $\gamma\cup\{j\}$ and $j$, and
    \cref{lemma:noise}\ref{it:noise_ii} applied to the nested pair
    $\gamma\subset\gamma\cup\{j\}$,
    \begin{equation*}
        \norm{\Pi Y}_2^{2}
        \ \leq\ 2\norm{\Pi\signal}_2^{2}+2\norm{\Pi\widetilde w}_2^{2}
        \ \leq\ 2n\,\omega(X)
        \bigl\|\beta^{\ast}_{\gamma^{\ast}\setminus\gamma}\bigr\|_2^{2}
        +4(L+\Ltil)\sigma_0^{2}\log p ,
    \end{equation*}
    and $\|\beta^{\ast}_{\gamma^{\ast}\setminus\gamma}\|_2^{2}
    \leq\|\beta^{\ast}_{\gamma^{\ast}}\|_2^{2}$. Multiplying by
    $4/\sigma_0^{2}$ gives~\eqref{eq:a0}.
\end{proof}

\subsection{Proofs for \cref{sec:flow}}
\label{sub:appflow}

\begin{proof}[Proof of \cref{lemma:funnel}]
    \ref{it:fn_i} None of \ref{H:null}--\ref{H:del} changes
    $\gamma\cap C_\lambda$, by the three sentences that prove
    \cref{lemma:Gvalid}\ref{it:gv_ii}, so $F_\lambda$ preserves $B_J$.
    The function~\eqref{eq:Vprog} without its fourth term drops by
    exactly one at each step of $F_\lambda$ and vanishes precisely when
    $\gamma\subseteq\gamma^{\ast}$, $c^{-}_\lambda\subseteq\gamma$ and
    $\gamma\subseteq c^{+}_\lambda$; within $B_J$ the only such state is
    $c^{-}_\lambda\cup J=v_J$. Hence every $\gamma\in B_J$ reaches $v_J$,
    which is~\ref{it:fn_i}.

    \ref{it:fn_ii} The three drop bounds of \cref{th:dropG} for rules
    \ref{H:null}--\ref{H:del} are all at most
    $\max\{p^{-\eta},(2s^{\ast})^{-1}\}$, which is at most $1$; the
    length and single-flip claims are
    \cref{lemma:Gvalid}\ref{it:gv_i} and $V\leq s_0+s^{\ast}\leq2s_0$.

    \ref{it:fn_iii} Run the proof of \cref{cor:precG} for $F_\lambda$ at
    the state $v_J$. Since rule~\ref{H:undec} is absent, the third factor
    is missing and the first two give
    $\pi_n^{(\lambda)}[\Lambda^{F_\lambda}(v_J)\mid Y]\leq
    2(2s^{\ast}+1)\,\pi_n^{(\lambda)}(v_J\mid Y)$. By~\ref{it:fn_i} the
    left-hand side is $\bar\pi_\lambda(J)$.
\end{proof}

\begin{proof}[Proof of \cref{th:gapflow}]
    Every edge of $\widetilde{\mathcal{T}}_\lambda$ is a single flip, so
    \cref{lemma:kernel} gives
    $\mathbf{Q}_\lambda(e)\geq\tfrac{1}{4p}\min\{\pi_n^{(\lambda)}
    (\bar\gamma\mid Y),\pi_n^{(\lambda)}(\bar\gamma'\mid Y)\}$ for its two
    endpoints. There are two kinds.

    \emph{Edges of segments (i) and (iii).} Such an edge is
    $e=(\bar\gamma,F_\lambda(\bar\gamma))$ for some $\bar\gamma$, and by
    \cref{lemma:funnel}\ref{it:fn_ii} it is uphill, so the minimum is
    $\pi_n^{(\lambda)}(\bar\gamma\mid Y)$. A pair
    $(\gamma,\gamma')$ uses $e$ in the direction of $F_\lambda$ only
    within its segment~(i), where the argument of
    \cref{lemma:paths}\ref{it:pp_iii} applies unchanged and forces
    $\gamma\in\Lambda^{F_\lambda}(\bar\gamma)$; the opposite orientation
    is the mirror image and gives the same bound because
    $\mathbf{Q}_\lambda$ is symmetric. Hence, using
    $\sum_{\gamma'}\pi_n^{(\lambda)}(\gamma'\mid Y)=1$ and
    \cref{lemma:funnel}\ref{it:fn_iii} applied at $\bar\gamma$ in place of
    $v_J$, that is the bound
    $\pi_n^{(\lambda)}[\Lambda^{F_\lambda}(\bar\gamma)\mid Y]\leq
    2(2s^{\ast}+1)\pi_n^{(\lambda)}(\bar\gamma\mid Y)$ established in its
    proof,
    \begin{equation*}
        \frac{1}{\mathbf{Q}_\lambda(e)}
        \sum_{\widetilde T_{\gamma,\gamma'}\ni e}
        \pi_n^{(\lambda)}(\gamma\mid Y)\pi_n^{(\lambda)}(\gamma'\mid Y)
        \ \leq\ 4p\,\frac{\pi_n^{(\lambda)}
            [\Lambda^{F_\lambda}(\bar\gamma)\mid Y]}
        {\pi_n^{(\lambda)}(\bar\gamma\mid Y)}
        \ \leq\ 8p\,(2s^{\ast}+1) .
    \end{equation*}

    \emph{Edges of segment (ii).} Such an edge is
    $e=(v_J,v_{J\triangle\{j\}})$ with $j\in C_\lambda$, and it is used by
    $(\gamma,\gamma')\in B_{J_1}\times B_{J_2}$ exactly when the
    increasing-order path from $J_1$ to $J_2$ traverses
    $(J,J\triangle\{j\})$. Summing $\pi_n^{(\lambda)}(\gamma\mid Y)
    \pi_n^{(\lambda)}(\gamma'\mid Y)$ over each block gives
    $\bar\pi_\lambda(J_1)\bar\pi_\lambda(J_2)$, so the numerator is the
    inner sum of~\eqref{eq:rholam}. For the denominator,
    \cref{lemma:funnel}\ref{it:fn_iii} gives
    \begin{equation*}
        \mathbf{Q}_\lambda(e)\ \geq\ \frac{1}{4p}
        \min\bigl\{\pi_n^{(\lambda)}(v_J\mid Y),
        \pi_n^{(\lambda)}(v_{J\triangle\{j\}}\mid Y)\bigr\}
        \ \geq\ \frac{\min\{\bar\pi_\lambda(J),
            \bar\pi_\lambda(J\triangle\{j\})\}}{8p\,(2s^{\ast}+1)} ,
    \end{equation*}
    so the congestion at $e$ is at most
    $8p(2s^{\ast}+1)\varrho_\lambda$.

    \emph{Length.} Segments (i) and (iii) have at most $2s_0$ edges each
    by \cref{lemma:funnel}\ref{it:fn_ii} and segment~(ii) at most
    $k_\lambda\leq s^{\ast}\leq s_0$, so
    $\ell(\widetilde{\mathcal{T}}_\lambda)\leq5s_0$. Sinclair's bound and
    \cref{lemma:kernel} now give~\eqref{eq:gapflow}.
\end{proof}

\begin{proof}[Proof of \cref{lemma:prodfree}]
    Write $\bar\pi_\lambda=\bigotimes_{i=1}^{k_\lambda}\mu_i$ and fix the
    directed edge $(J,J\triangle\{j\})$. A pair $(J_1,J_2)$ traverses it
    exactly when, at that moment, the coordinates below $j$ have been
    flipped and those above have not: that is, $J_2$ agrees with $J$
    below $j$ and takes the value $1-J_j$ at $j$, while $J_1$ agrees with
    $J$ at $j$ and above. So $J_1$ is free below $j$ and $J_2$ is free
    above $j$, and by independence
    \begin{multline*}
        \sum_{(J_1,J_2)}\bar\pi_\lambda(J_1)\bar\pi_\lambda(J_2)
        =\Bigl[\sum_{a}\textstyle\prod_{i<j}\mu_i(a_i)\Bigr]
        \prod_{i\geq j}\mu_i(J_i)\\
        \cdot\ \prod_{i<j}\mu_i(J_i)\,\mu_j(1-J_j)
        \Bigl[\sum_{b}\textstyle\prod_{i>j}\mu_i(b_i)\Bigr] ,
    \end{multline*}
    the two bracketed sums being $1$. The remaining product is
    $\bar\pi_\lambda(J)\,\mu_j(1-J_j)$. Since
    $\bar\pi_\lambda(J\triangle\{j\})=\bar\pi_\lambda(J)\mu_j(1-J_j)
    /\mu_j(J_j)$, the quotient in~\eqref{eq:rholam} equals
    \begin{equation*}
        \frac{\bar\pi_\lambda(J)\,\mu_j(1-J_j)}
        {\bar\pi_\lambda(J)\min\bigl\{1,\mu_j(1-J_j)/\mu_j(J_j)\bigr\}}
        =\max\bigl\{\mu_j(J_j),\,\mu_j(1-J_j)\bigr\}\ \leq\ 1 . \qedhere
    \end{equation*}
\end{proof}

\begin{proof}[Proof of \cref{lemma:comparison}]
    The inner sum of~\eqref{eq:rholam} is over the same index set for
    $\bar\pi_\lambda$ and $\bar\mu$, and each of its terms is a product of
    two masses, so it is at most $e^{2D}$ times the corresponding sum for
    $\bar\mu$; the denominator is at least $e^{-D}$ times the
    corresponding minimum. Apply \cref{lemma:prodfree} to $\bar\mu$.
\end{proof}

\begin{proof}[Proof of \cref{prop:rhobound}]
    Let $\bar\mu$ be the product measure on $2^{C_\lambda}$ with odds
    $\theta_j\coloneq\exp\{\lambda\psi_j(c^{-}_\lambda)-\kappa\log p\}$,
    normalised, so that
    $\bar\mu(J)\propto\prod_{j\in J}\theta_j$. Enumerating
    $J=\{j_1<\dots<j_m\}$ and telescoping~\eqref{eq:onestepG} along
    $c^{-}_\lambda\subset c^{-}_\lambda\cup\{j_1\}\subset\dots\subset v_J$,
    \begin{equation*}
        \log\frac{\pi_n^{(\lambda)}(v_J\mid Y)}
        {\pi_n^{(\lambda)}(v_\emptyset\mid Y)}
        =\sum_{i=1}^{m}\bigl\{\lambda\psi_{j_i}
        (c^{-}_\lambda\cup\{j_1,\dots,j_{i-1}\})-\kappa\log p\bigr\},
    \end{equation*}
    and each summand differs from $\log\theta_{j_i}$ by at most
    $\lambda\Delta^{C}_\lambda\leq\Delta^{C}_\lambda$
    by~\eqref{eq:DeltaC}, the contexts involved being of the form
    $c^{-}_\lambda\cup U$ with $U\subseteq C_\lambda\setminus\{j_i\}$.
    Writing $\widetilde\mu(J)\coloneq\prod_{j\in J}\theta_j$ and
    $\widetilde g(J)\coloneq\pi_n^{(\lambda)}(v_J\mid Y)/
    \pi_n^{(\lambda)}(v_\emptyset\mid Y)$ for the two unnormalised
    weights, the display therefore gives
    \begin{equation}
        \label{eq:D1}
        \Bigl|\log\frac{\widetilde g(J)}{\widetilde\mu(J)}\Bigr|
        \ \leq\ \abs{J}\,\Delta^{C}_\lambda
        \ \leq\ k_\lambda\Delta^{C}_\lambda
        \ \eqcolon\ D_1
        \qquad\text{for every }J\subseteq C_\lambda .
    \end{equation}
    Summing~\eqref{eq:D1} over $J$, the two normalising constants
    $\sum_J\widetilde g$ and $\sum_J\widetilde\mu=Z_\mu$ also differ by a
    factor at most $e^{D_1}$, so the two \emph{normalised} measures
    differ by a factor at most $e^{2D_1}$ pointwise.

    Next, \cref{lemma:funnel}\ref{it:fn_iii} and
    $\pi_n^{(\lambda)}(v_J\mid Y)\leq\bar\pi_\lambda(J)$ give
    $\bar\pi_\lambda(J)=\pi_n^{(\lambda)}(v_J\mid Y)\,c(J)$ with
    $c(J)\in[1,2(2s^{\ast}+1)]$; both $\bar\pi_\lambda$ and the
    normalisation of $\widetilde g$ being probability measures, they
    differ pointwise by a factor at most $\{2(2s^{\ast}+1)\}^{2}$, one
    factor from $c(J)$ and one from the normaliser.

    Combining the two paragraphs, \cref{lemma:comparison} applies with
    $D=2k_\lambda\Delta^{C}_\lambda+2\log\{2(2s^{\ast}+1)\}$, and
    $\varrho_\lambda\leq e^{3D}$ is exactly~\eqref{eq:rhobound}.
\end{proof}
